\documentclass[11pt]{article}

\usepackage[margin=1in]{geometry}
\usepackage{amsmath}
\usepackage{booktabs}
\usepackage{caption}
\usepackage{graphicx}
\usepackage{hyperref}
\usepackage{pdflscape}
\usepackage{array}

\hypersetup{
  colorlinks=true,
  linkcolor=blue,
  urlcolor=blue,
  citecolor=blue
}

\graphicspath{{../results/}}

\title{IFEval Factor Analysis: Rank Selection, Sparsity Tuning, Predictive Comparison, and Selected-Model Interpretation}
\author{}
\date{\today}

\begin{document}
\maketitle

\section{Overview}

This note summarizes the reproducible IFEval factor-analysis pipeline. The data are a binary
model-by-item response matrix with 122 model rows and 500 IFEval item columns. The analysis has
three goals:
\begin{enumerate}
  \item select the working factor rank using held-out predictive performance;
  \item tune the entrywise sparsity penalty on the loading matrix;
  \item compare the selected independent-mixture probit factor model to an ordinary binary probit
        factor baseline;
  \item interpret and visualize the selected factor space through loadings, cross-loadings,
        mixture profiles, and three-dimensional factor scores.
\end{enumerate}

\section{Models}

The selected model is an independent marginal-mixture probit factor model,
\begin{align}
X_{ij} \mid f_i, \alpha_j, \lambda_j
  &\sim \mathrm{Bernoulli}\{\Phi(\alpha_j + \lambda_j^\top f_i)\}, \\
f_{ih}
  &\sim \sum_{g=1}^{G} \pi_{hg} N(\mu_{hg}, \sigma_{hg}^2),
  \qquad h=1,\ldots,H,
\end{align}
with independent mixture coordinates across factors. For the selected IFEval analysis, the
working model uses $H=3$ factors and $G=3$ marginal mixture components.

The main comparison model is an ordinary binary probit factor model,
\begin{align}
X_{ij} \mid f_i, \alpha_j, \lambda_j
  &\sim \mathrm{Bernoulli}\{\Phi(\alpha_j + \lambda_j^\top f_i)\}, \\
f_i &\sim N(0, I_H),
\end{align}
also fit at $H=3$ for the selected-model comparison. Because factor labels and signs are
unidentified, the ordinary probit factors are aligned to the mixture factors by the signed
permutation that maximizes factor-score correlations.

\section{Rank Selection and Predictive Evaluation}

Rank selection is based on held-out cell predictive log likelihood over $H=1,\ldots,10$. The
comparison includes independent-mixture probit fits with $G=2$ and $G=3$, along with the ordinary
Gaussian-factor probit baseline. All three methods select $H=3$ by the best mean held-out log
score, and the one-standard-error rule also selects $H=3$.

\begin{table}[h]
\centering
\caption{Selected rank and held-out log score by method.}
\label{tab:rank-selection}
\begin{tabular}{lrrrr}
\toprule
Method & Best $H$ & Mean log score & SE & Mean seconds at best $H$ \\
\midrule
Mixture probit, $G=2$ & 3 & -0.2762 & 0.0039 & 3.285 \\
Mixture probit, $G=3$ & 3 & -0.2756 & 0.0033 & 3.434 \\
Ordinary probit factor & 3 & -0.2792 & 0.0041 & 3.062 \\
\bottomrule
\end{tabular}
\end{table}

The $G=3$ mixture has the strongest held-out log score, though the differences are modest relative
to the fold-to-fold uncertainty. The main practical conclusion is therefore not only that $H=3$
is selected, but also that the mixture model provides a slightly better predictive account without
changing the selected rank.

\begin{table}[h]
\centering
\caption{Held-out predictive metrics at $H=3$.}
\label{tab:predictive-metrics}
\begin{tabular}{lrrrrrr}
\toprule
Method & Accuracy & SE & Brier & SE & Log score & SE \\
\midrule
Mixture $G=2$ & 0.8861 & 0.0009 & 0.0824 & 0.0010 & -0.2762 & 0.0039 \\
Mixture $G=3$ & 0.8866 & 0.0008 & 0.0822 & 0.0008 & -0.2756 & 0.0033 \\
Ordinary probit & 0.8852 & 0.0017 & 0.0826 & 0.0011 & -0.2792 & 0.0041 \\
\bottomrule
\end{tabular}
\end{table}

\begin{figure}[h]
\centering
\includegraphics[width=0.86\textwidth]{reproduced_openeval_ifeval_rank_selection_comparison/openeval_ifeval_heldout_loglik_by_H_method.png}
\caption{Held-out predictive log likelihood across candidate ranks.}
\label{fig:rank-curve}
\end{figure}

\begin{figure}[h]
\centering
\includegraphics[width=0.94\textwidth]{../writeup/ifeval_H3_predictive_metrics_clean.png}
\caption{Held-out predictive metrics for the selected $H=3$ fits.}
\label{fig:predictive-metrics}
\end{figure}

\section{Lambda Sparsity Tuning}

After selecting $H=3$ and $G=3$, the entrywise L1 penalty on the loading matrix is tuned over the
grid $\{0,0.5,1,2,3,4,6,8,10,12\}$. The tuning path uses a single common pretraining fit and then
reruns the refinement stage for each candidate value. We summarize the tradeoff between unpenalized
binary probit log likelihood and loading sparsity, and select the penalty by a BIC-style criterion
in which the varying degrees of freedom are driven by the number of nonzero loading entries.

\begin{table}[h]
\centering
\caption{Entrywise loading-penalty tuning path for the selected $H=3,G=3$ model.}
\label{tab:lambda-tuning}
\begin{tabular}{rrrrr}
\toprule
$\lambda_{\mathrm{L1}}$ & Log score & Nonzero loadings & Density & BIC-style score \\
\midrule
0.0 & -0.2166 & 1500 & 1.000 & 52762.9 \\
0.5 & -0.2185 & 1447 & 0.965 & 52402.0 \\
1.0 & -0.2202 & 1406 & 0.937 & 52158.0 \\
2.0 & -0.2237 & 1329 & 0.886 & 51740.0 \\
3.0 & -0.2274 & 1262 & 0.841 & 51455.8 \\
4.0 & -0.2313 & 1198 & 0.799 & 51223.7 \\
6.0 & -0.2403 & 967 & 0.645 & 49780.5 \\
8.0 & -0.2479 & 837 & 0.558 & 49276.4 \\
10.0 & -0.2549 & 759 & 0.506 & 49268.5 \\
12.0 & -0.2617 & 717 & 0.478 & 49628.4 \\
\bottomrule
\end{tabular}
\end{table}

The raw likelihood favors less shrinkage, as expected, but the BIC-style criterion favors
$\lambda_{\mathrm{L1}}=10$. The criterion decreases through $\lambda_{\mathrm{L1}}=8$, reaches
its minimum at 10, and then increases at 12, so the expanded grid reaches the tuning shelf rather
than selecting a boundary value. The selected penalty keeps about half of the loading entries
nonzero in the tuning path and reduces the number of large entries with $|\lambda|>0.5$ to 413.

\begin{figure}[h]
\centering
\includegraphics[width=0.92\textwidth]{reproduced_openeval_ifeval_lambda_sparsity_tuning/ifeval_lambda_sparsity_tuning_path.png}
\caption{Loading sparsity tuning path. The dashed line marks the BIC-selected penalty.}
\label{fig:lambda-tuning-path}
\end{figure}

\section{Selected-Model Fit}

The selected independent-mixture fit uses $H=3$, $G=3$, and the tuned loading penalty
$\lambda_{\mathrm{L1}}=10$. The fitted loading matrix has a clear
asymmetry across axes: factor 1 is broad and dense, while factors 2 and 3 are more selective.
Using a threshold $|\lambda_{jh}| \geq 0.5$, 349 of 500 items load on factor 1, compared with 30
items on factor 2 and 34 items on factor 3. This pattern supports reading factor 1 as a broad
instruction-following or overall compliance axis, with factors 2 and 3 capturing more specialized
modes of instruction control. The reported sparse loading matrix sets all entries below this
threshold to zero and is saved as
\path{results/loadings_crossloadings/ifeval_sparse_lambda_matrix_threshold_0p5.csv}.

\begin{table}[h]
\centering
\caption{Summary of fitted loading magnitudes.}
\label{tab:lambda-summary}
\begin{tabular}{lrrrr}
\toprule
Factor & Mean $|\lambda|$ & Median & $|\lambda|>0.5$ & $|\lambda|>1.0$ \\
\midrule
F1 & 0.665 & 0.710 & 349 & 94 \\
F2 & 0.084 & 0.000 & 30 & 24 \\
F3 & 0.130 & 0.000 & 34 & 5 \\
\bottomrule
\end{tabular}
\end{table}

\begin{landscape}
\begin{figure}[p]
\centering
\includegraphics[width=0.90\linewidth,height=0.78\textheight,keepaspectratio]{reproduced_openeval_ifeval_H3_G3_interpretation/openeval_lambda_full_ordered_by_benchmark.png}
\caption{Selected mixture-model loading matrix, ordered by benchmark.}
\label{fig:mixture-loadings-benchmark}
\end{figure}
\end{landscape}

\section{Comparison to Ordinary Binary Probit}

The ordinary binary probit baseline is fit at the same selected rank, $H=3$. After signed
permutation alignment, its factor space is compared directly with the mixture factor space. This
comparison asks whether the predictive and interpretive gains from the mixture model are merely a
rotation of a Gaussian factor model, or whether the mixture prior produces a more structured
profile geometry.

\begin{figure}[h]
\centering
\includegraphics[width=\textwidth,height=0.58\textheight,keepaspectratio]{reproduced_openeval_ordinary_vs_mixture_H3/openeval_lambda_side_by_side_ordered_by_benchmark.png}
\caption{Aligned loading comparison: independent-mixture probit versus ordinary binary probit.}
\label{fig:lambda-side-by-side}
\end{figure}

The side-by-side loading plot shows that the broad first-factor structure is shared to some extent,
but the mixture model produces a more interpretable partition into marginal mixture profiles. That
profile structure is especially useful when reading the factor scores as model types rather than
only as continuous latent coordinates.

\section{Factor Interpretation}

All of the top-loading item sets are IFEval instruction-following prompts, so the factor
interpretation is about modes of instruction control rather than benchmark-domain differences.
The positive side of factor 1 contains many direct compliance constraints, such as required
wrapping, exact ending phrases, paragraph counts, JSON format, capitalization, and no-comma
constraints. This factor behaves like a broad general instruction-following axis.

Factors 2 and 3 appear more specialized. Factor 2 has a small number of very large loadings,
including prompts requiring strict casing, language restrictions, list or poem formats, and
multi-constraint style transformations. Factor 3 loads on a larger but still selective set of
items. In prior inspection, these items looked especially connected to creative or generative
control: producing songs, poems, essays, stories, or styled rewrites while also respecting formal
constraints. The model therefore separates broad compliance from more specific forms of stylistic
and creative control.

\begin{table}[h]
\centering
\caption{Top-loading item summaries by factor and direction.}
\label{tab:factor-summary}
\begin{tabular}{llrrr}
\toprule
Factor & Direction & Mean signed loading & Mean $|\lambda|$ & Mean item accuracy \\
\midrule
F1 & Positive & 1.196 & 1.196 & 0.737 \\
F1 & Negative & 0.051 & 0.054 & 0.600 \\
F2 & Positive & 0.559 & 0.559 & 0.551 \\
F2 & Negative & -0.004 & 0.004 & 0.773 \\
F3 & Positive & 0.527 & 0.527 & 0.693 \\
F3 & Negative & -0.040 & 0.040 & 0.777 \\
\bottomrule
\end{tabular}
\end{table}

\section{Cross-Loading Structure}

The thresholded loading analysis gives 375 factor-only items and 19 cross-loading items at
$|\lambda| \geq 0.5$. Among factor-only items, 330 are strongest on F1, 30 on F2, and 15 on F3.
All cross-loading items involve F1 and F3, with active pattern 101. This reinforces the view
that F1 is broad, while F2 and F3 capture refinements of instruction control that often appear
jointly with the broad axis.

\begin{figure}[h]
\centering
\includegraphics[width=0.68\textwidth]{loadings_crossloadings/ifeval_G3_cross_loading_items_heatmap.png}
\caption{Cross-loading items under the threshold $|\lambda| \geq 0.5$.}
\label{fig:cross-loadings}
\end{figure}

\section{Mixture Profiles and Factor Visualization}

The three marginal mixture coordinates induce a profile label for each model. The highest-accuracy
profile is 3-3-3, with 12 models and mean accuracy 0.959. Other high-performing groups include
3-3-2, with 16 models and mean accuracy 0.920, and 3-1-3, with 20 models and mean accuracy 0.898.
The profile labels therefore provide a compact summary of model capability across the broad
instruction-following axis and the two more specialized axes.

\begin{figure}[h]
\centering
\includegraphics[width=0.88\textwidth]{ifeval_3d_factor_visualizations/ifeval_3d_factor_scatter_profiles.png}
\caption{Selected mixture-model factor scores in three dimensions, colored by MAP mixture profile.}
\label{fig:mixture-3d}
\end{figure}

\begin{figure}[h]
\centering
\includegraphics[width=0.94\textwidth]{reproduced_ifeval_ordinary_vs_mixture_factor_visualization/ifeval_mixture_vs_ordinary_probit_3d_profiles.png}
\caption{Static three-dimensional projection comparing mixture factors to aligned ordinary probit factors.}
\label{fig:ordinary-vs-mixture-3d}
\end{figure}

\clearpage

\section{Summary}

The full IFEval pipeline selects a three-factor representation by held-out predictive evaluation,
then selects an entrywise loading penalty $\lambda_{\mathrm{L1}}=10$ by a BIC-style sparsity
tuning criterion. At $H=3$, the $G=3$ independent-mixture probit fit has the best held-out log
score and Brier score among the compared methods, while the ordinary probit factor baseline is
close but slightly worse.
The selected loading matrix is interpretable: F1 is a broad instruction-following factor, while
F2 and F3 capture more specialized modes of constraint satisfaction and creative/stylistic
control. The cross-loading analysis shows that many specialized items still involve the broad
axis, especially the F1/F3 pattern, which is consistent with a block-like IFEval structure rather
than a strict bifactor model. The three-dimensional factor visualization makes this profile
structure visible at the model level and gives a useful bridge between predictive evaluation and
qualitative interpretation.

\end{document}
